\chapter{Wittig y otras reacciones de formación de C=C}\label{ch:b2:wittig}

La hembra de la mosca doméstica atrae a los machos con un hidrocarburo de 23 carbonos que lleva un solo doble enlace, entre
los carbonos 9 y 10, de geometría (\textit{Z}). Para fabricarlo en cantidad, el químico debe colocar ese único
doble enlace exactamente ahí, y con esa geometría: una eliminación lo dispersaría por varias
posiciones y daría las dos geometrías. Este capítulo construye los dobles enlaces por el otro camino, uniendo dos
fragmentos por un grupo carbonilo, de modo que el doble enlace aparece donde estaba el \ce{C=O}.
% ledger: use:muscalure

\begin{recall}
\Cref{ch:b2:enolates-aldol}: la \omterm{def:b2:enolates-aldol:aldol}{condensación aldólica}. \Cref{ch:b2:alkene-redox}: la hidrogenación de
un alquino a (\textit{Z})-alqueno sobre un catalizador envenenado. El volumen del primer año: E1, E2 y regla de
Zaitsev, sustitución bimolecular, reactivos de Grignard e inversión de polaridad, descriptores
(\textit{Z})/(\textit{E}). El volumen escolar: economía atómica.
\end{recall}

\section{Sales de fosfonio e iluros}

\begin{definition}[Iluros de fosfonio]\label{def:b2:wittig:ylide}
Una \emph{sal de fosfonio}\index{sal de fosfonio} $\mathrm{R_3P^+{-}CH_2R'\ X^-}$ se prepara por sustitución
bimolecular de un halogenoalcano con una fosfina, normalmente la trifenilfosfina \ce{Ph3P}. Arrancar un
hidrógeno del carbono unido al fósforo da un \emph{iluro de fosfonio}\index{iluro de fosfonio},
una molécula neutra con cargas opuestas contiguas, $\mathrm{Ph_3P^+{-}CH^-{-}R'}$, que también se escribe $\mathrm{Ph_3P{=}CHR'}$. Un
\emph{iluro estabilizado}\index{iluro estabilizado} lleva sobre ese carbono un grupo atractor de electrones
(un éster, una cetona, un \omterm{def:b2:amines:nitrile}{nitrilo}) que deslocaliza la carga negativa.
\end{definition}

\begin{omfigure}
\setchemfig{atom sep=1.9em}
\schemestart
\chemfig{Ph_3P} \+ \chemfig{CH_3-I}
\arrow{->[S$_\mathrm{N}$2]}[,0.9]
\chemfig{Ph_3P^{+}-CH_3}\;\chemfig{I^{-}}
\arrow{->[\ce{BuLi}][$-$ \ce{BuH}, \ce{LiI}]}[,1.3]
\chemfig{Ph_3P^{+}-\chemabove{C}{\scriptstyle -}H_2}
\arrow{<->}[,0.8]
\chemfig{Ph_3P=CH_2}
\schemestop
\par\vspace{6pt}

{\small El metilidentrifenilfosforano, el iluro más sencillo. La trifenilfosfina desplaza el yoduro
del yodometano; el butil-litio arranca un protón del grupo metilo. Las dos estructuras de la derecha
son dos maneras de escribir la misma molécula: la de cargas separadas es la más cercana a la realidad.}
\end{omfigure}

\begin{proposition}[Preparar el iluro]\label{prop:b2:wittig:ylide-acidity}
Una \omterm{def:b2:wittig:ylide}{sal de fosfonio} es mucho más ácida en el carbono contiguo al fósforo que un alcano. Una sal de alquilfosfonio
sencilla necesita aun así una base muy fuerte (butil-litio, hidruro de sodio, amiduro de sodio) en condiciones anhidras
e inertes; una sal cuyo \ce{CH2} lleva además un grupo carbonilo se desprotona con una base débil
(carbonato de sodio, hidróxido de sodio), y el \omterm{def:b2:wittig:ylide}{iluro estabilizado} formado puede aislarse y almacenarse.
\end{proposition}

\begin{proof}[Argumento]
El fósforo positivo contiguo al carbanión lo estabiliza por su carga y por la polarizabilidad del
voluminoso átomo de fósforo: un iluro sencillo es un carbanión estabilizado, pero menos que un \omterm{def:b2:enolates-aldol:enolate}{enolato},
de ahí la base fuerte. Un carbonilo contiguo reparte la carga negativa sobre el oxígeno, como en un \omterm{def:b2:enolates-aldol:enolate}{enolato}
(\cref{prop:b2:enolates-aldol:alpha-acidity}), además de sobre el fósforo: el ácido conjugado se vuelve
lo bastante ácido para una base débil, y el iluro, mucho menos básico, ya no reacciona con el agua ni con el aire.
\end{proof}

\section{La reacción de Wittig}

\begin{definition}[Reacción de Wittig]\label{def:b2:wittig:wittig}
La \emph{reacción de Wittig}\index{reaccion de Wittig@reacción de Wittig} de un \omterm{def:b2:wittig:ylide}{iluro de fosfonio} con un aldehído o una cetona
da un alqueno y óxido de trifenilfosfina:
\begin{equation*}
\mathrm{Ph_3P{=}CHR + R'CHO \longrightarrow R'CH{=}CHR + Ph_3P{=}O}.
\end{equation*}
Pasa por un \emph{oxafosfetano}\index{oxafosfetano}, un anillo de cuatro eslabones formado por el fósforo, dos
carbonos y el oxígeno, que se forma cuando el carbono del iluro se une al carbono carbonílico mientras el oxígeno carbonílico
se une al fósforo.
\end{definition}

\begin{omfigure}
\setchemfig{atom sep=1.9em}
\schemestart
\chemfig{Ph_3P=CH-R} \+ \chemfig{O=CH-R'}
\arrow{->[adición]}[,1.0]
\chemfig{Ph_3P?-[6]CH(-[4]R)-[0]CH(-[0]R')-[2]O?}
\schemestop

\medskip
\schemestart
\arrow{->[eliminación]}[,1.1]
\chemfig{R-CH=CH-R'} \+ \chemfig{Ph_3P=O}
\schemestop
\par\vspace{6pt}

{\small La \omterm{def:b2:wittig:wittig}{reacción de Wittig}. El carbono del iluro y el carbono carbonílico se unen mientras el oxígeno se une
al fósforo, en una sola etapa: el \omterm{def:b2:wittig:wittig}{oxafosfetano}. Su anillo se abre después en el otro sentido, rompiendo los
enlaces \ce{P-C} y \ce{C-O}: el enlace \ce{C=C} se forma entre los dos antiguos carbonos reaccionantes, y el
enlace \ce{P=O} del óxido de trifenilfosfina se forma al mismo tiempo.}
\end{omfigure}

\begin{proposition}[Un doble enlace exactamente en su sitio]\label{prop:b2:wittig:regiospecific}
En una \omterm{def:b2:wittig:wittig}{reacción de Wittig} el nuevo enlace \ce{C=C} une el antiguo carbono carbonílico al antiguo carbono del iluro,
y a ningún otro.
\end{proposition}

\begin{proof}[Argumento]
Los únicos enlaces que se forman y se rompen son los del anillo de cuatro eslabones: el carbono del iluro con el carbono
carbonílico, y luego los enlaces \ce{C-P} y \ce{C-O}. Ningún hidrógeno se desplaza ni se forma carbocatión ni carbanión
a lo largo de la cadena, de modo que el doble enlace no puede migrar. Una eliminación, en cambio, elige entre los
hidrógenos de varios carbonos vecinos (regla de Zaitsev), y una reacción E1 puede sufrir transposiciones.
\end{proof}

\begin{proposition}[Fuerza motriz]\label{prop:b2:wittig:driving}
La formación del enlace \ce{P=O} del óxido de trifenilfosfina hace que la \omterm{def:b2:wittig:wittig}{reacción de Wittig} sea muy
favorable e irreversible.
\end{proposition}

\begin{proof}[Argumento]
La reacción cambia un enlace \ce{C=O} y uno \ce{P-C} (en el iluro, el \ce{P=C} es sobre todo un enlace sencillo
\ce{P+-C-}) por un enlace \ce{C=C} y uno \ce{P=O}. El enlace \ce{P=O} de un óxido de fosfina es
muy fuerte, mucho más que el enlace \ce{P-C} perdido, y el \ce{C=O} perdido queda compensado en parte
por el \ce{C=C} formado: el balance es claramente favorable. El fósforo es el sumidero de oxígeno de la
reacción.
\end{proof}

\section{Estereoselectividad}

\begin{proposition}[Geometría del alqueno]\label{prop:b2:wittig:stereo}
Con un aldehído, un iluro no estabilizado (\ce{Ph3P=CHR}, R un grupo alquilo) da principalmente el (\textit{Z})-alqueno,
en condiciones exentas de sales. Un \omterm{def:b2:wittig:ylide}{iluro estabilizado} (\ce{Ph3P=CH-COOR}, \ce{Ph3P=CH-COR}) da principalmente el
(\textit{E})-alqueno. Un \omterm{def:b2:wittig:ylide}{iluro estabilizado} solo por un grupo arilo da mezclas.
\end{proposition}

\admitted

Son observaciones. Su explicación habitual lee la \omterm{def:b2:wittig:wittig}{reacción de Wittig} como una competencia entre
velocidades y estabilidades. Con un iluro reactivo, no estabilizado, el \omterm{def:b2:wittig:wittig}{oxafosfetano} se forma deprisa y de forma
irreversible, a través del estado de transición menos congestionado, en el que los dos sustituyentes acaban del
mismo lado del anillo; después colapsa conservando esa disposición: el control cinético da el
(\textit{Z})-alqueno. Con un \omterm{def:b2:wittig:ylide}{iluro estabilizado}, la primera etapa es más lenta y reversible; los
\omterm{def:b2:wittig:wittig}{oxafosfetanos} se interconvierten a través de los productos de partida, y gana el que tiene los sustituyentes en
lados opuestos, menos tensionado: el (\textit{E})-alqueno, bajo control termodinámico.

\begin{omfigure}
\setchemfig{atom sep=1.9em}
\schemestart
\chemfig{Ph_3P=CH-CH_2CH_3} \+ \chemfig{O=CH-Ph}
\arrow{->}[,0.8]
\chemfig{Ph-[:-30]=[:30]-[:90]CH_2CH_3}
\schemestop

\medskip
\schemestart
\chemfig{Ph_3P=CH-COOEt} \+ \chemfig{O=CH-Ph}
\arrow{->}[,0.8]
\chemfig{Ph-[:-30]=[:30]-[:-30]COOEt}
\schemestop
\par\vspace{6pt}

{\small Dos iluros, un aldehído. Arriba: el iluro de propilideno, no estabilizado, da principalmente
(\textit{Z})-1-fenilbut-1-eno. Abajo: el \omterm{def:b2:wittig:ylide}{iluro estabilizado} da principalmente
(\textit{E})-3-fenilpropenoato de etilo (cinamato de etilo). El óxido de trifenilfosfina, formado en ambos casos, no se
muestra.}
\end{omfigure}

\section{La reacción de Horner--Wadsworth--Emmons}

\begin{definition}[Fosfonatos y reacción HWE]\label{def:b2:wittig:hwe}
Un \emph{fosfonato}\index{fosfonato} es un éster de un ácido fosfónico, $\mathrm{(RO)_2P({=}O){-}R'}$, con un
enlace \ce{P-C}. Se prepara por la reacción de Michaelis--Arbuzov de un fosfito de trialquilo con un
halogenoalcano, $\mathrm{P(OEt)_3 + R'Br \to (EtO)_2P(O)R' + EtBr}$. En la
\emph{reacción de Horner--Wadsworth--Emmons}\index{reaccion de Horner--Wadsworth--Emmons@reacción de Horner--Wadsworth--Emmons} (HWE), el anión de un
fosfonato que lleva un grupo atractor de electrones en su \ce{CH2} se adiciona a un aldehído y da un
alqueno y un anión fosfato de dialquilo.
\end{definition}

\begin{omfigure}
\setchemfig{atom sep=1.8em}
\schemestart
\chemfig{EtO-P(=[2]O)(-[6]OEt)-CH_2-COOEt}
\arrow{->[\ce{NaH}][$-$ \ce{H2}]}[,0.9]
\chemfig{EtO-P(=[2]O)(-[6]OEt)-\chemabove{C}{\scriptstyle -}H-COOEt}
\schemestop

\medskip
\schemestart
\arrow{->[\ce{PhCHO}]}[,0.9]
\chemfig{*6(=-=(-[:30]=[:-30]-[:30](=[:90]O)-[:-30]O-[:30]Et)-=-)}
\+
\chemfig{EtO-P(=[2]O)(-[6]O^{-})-OEt}
\schemestop
\par\vspace{6pt}

{\small La \omterm{def:b2:wittig:hwe}{reacción de Horner--Wadsworth--Emmons} del fosfonoacetato de trietilo con el benzaldehído. El hidruro
de sodio arranca el hidrógeno situado entre el \omterm{def:b2:wittig:hwe}{fosfonato} y el éster; el anión se adiciona al aldehído,
y el intermedio de cuatro eslabones colapsa como en la \omterm{def:b2:wittig:wittig}{reacción de Wittig}. El producto es el
(\textit{E})-cinamato de etilo; el subproducto es el anión fosfato de dietilo, soluble en agua.}
\end{omfigure}

\begin{proposition}[Por qué gusta a los químicos la reacción HWE]\label{prop:b2:wittig:hwe}
La reacción HWE da principalmente el (\textit{E})-alqueno, y su subproducto fosforado, una sal soluble en agua,
se elimina con un simple lavado acuoso; además, el anión \omterm{def:b2:wittig:hwe}{fosfonato} es más nucleófilo que el
\omterm{def:b2:wittig:ylide}{iluro estabilizado} correspondiente y reacciona también con las cetonas.
\end{proposition}

\begin{proof}[Argumento]
El anión \omterm{def:b2:wittig:hwe}{fosfonato} es un carbanión estabilizado, de modo que su adición es reversible y la reacción transcurre
bajo control termodinámico: (\textit{E}), como con los \omterm{def:b2:wittig:ylide}{iluros estabilizados}. A diferencia de un \omterm{def:b2:wittig:ylide}{iluro de fosfonio},
lleva una carga negativa completa, no compensada por un catión vecino: es más reactivo. El
subproducto, \ce{(EtO)2PO2-}, es una sal iónica, mientras que el óxido de trifenilfosfina es un sólido orgánico
neutro que hay que separar del producto por cromatografía o cristalización.
\end{proof}

\section{Elegir un método}

\begin{method}[Desconexión de Wittig]\label{met:b2:wittig:wittig-disconnection}
\begin{enumerate}
  \item Córtese el enlace \ce{C=C} objetivo; un carbono pasa a ser un carbono carbonílico, y el otro, el carbono del iluro
        (una \omterm{def:b2:wittig:ylide}{sal de fosfonio}, a partir de un haluro).
  \item De las dos maneras de hacerlo, prefiérase aquella cuyo haluro es primario (sustitución bimolecular por
        \ce{Ph3P}) y cuyo compuesto carbonílico es un aldehído.
  \item Elíjase el reactivo según la geometría buscada: un iluro no estabilizado para (\textit{Z}), un
        \omterm{def:b2:wittig:ylide}{iluro estabilizado} o un \omterm{def:b2:wittig:hwe}{fosfonato} (HWE) para (\textit{E}).
  \item Compruébese que nada más en la molécula reacciona con la base utilizada.
\end{enumerate}
\end{method}

\begin{method}[Elegir una reacción de formación de C=C]\label{met:b2:wittig:choose-cc}
\begin{center}
\small
\begin{tabular}{@{}lll@{}}
\hline
Método & Posición del \ce{C=C} & Geometría \\
\hline
Eliminación E2 o E1 & regla de Zaitsev, mezclas & sobre todo (\textit{E}), mezclas \\
\omterm{def:b2:enolates-aldol:aldol}{Condensación aldólica} & conjugado con \ce{C=O} & (\textit{E}) \\
Wittig, iluro no estabilizado & en el antiguo \ce{C=O} & (\textit{Z}) \\
Wittig, \omterm{def:b2:wittig:ylide}{iluro estabilizado}; HWE & en el antiguo \ce{C=O} & (\textit{E}) \\
Alquino y luego catalizador envenenado & en el antiguo triple enlace & (\textit{Z}) \\
\hline
\end{tabular}
\end{center}
\medskip
Para unir dos fragmentos de tamaño parecido por un \ce{C=C}, la \omterm{def:b2:wittig:wittig}{reacción de Wittig} y sus parientes suelen ser
la primera elección; la metátesis de olefinas, que intercambia los extremos de dos alquenos, se verá en el volumen
del tercer año.
\end{method}

\begin{history}[El iluro que se convirtió en reactivo]
Georg Wittig, que estudiaba los iluros de fósforo a principios de la década de 1950, descubrió que el metilidentrifenilfosforano
convierte la benzofenona en 1,1-difenileteno y óxido de trifenilfosfina; con su alumno Ulrich
Schöllkopf convirtió la observación en un método general en 1954. La industria lo adoptó con rapidez, entre otras cosas para la
vitamina A, y Wittig compartió el Premio Nobel de Química de 1979.
\end{history}

\begin{safety}{\ghs[5mm]{GHS02}\,\ghs[5mm]{GHS05}\,\ghs[5mm]{GHS07}\,\ghs[5mm]{GHS08}}
Las disoluciones de butil-litio son inflamables y corrosivas y reaccionan violentamente con el agua; el hidruro de sodio
desprende hidrógeno con el agua. La trifenilfosfina es nociva y corrosiva para los ojos, el fosfito de
trietilo, inflamable y nocivo; el óxido de trifenilfosfina es nocivo.
% ledger: ghs:BuLi, ghs:PPh3, ghs:triethyl-phosphite, ghs:Ph3PO
\end{safety}

\section{Ejercicios}

\begin{exercise}[$\star$]\label{exo:b2:wittig:1}
Escríbase la preparación del iluro \ce{Ph3P=CHCH3} a partir de la trifenilfosfina y un halogenoalcano.
\end{exercise}

\begin{exercise}[$\star$]\label{exo:b2:wittig:2}
Dense los productos de la \omterm{def:b2:wittig:wittig}{reacción de Wittig} del benzaldehído con el iluro preparado a partir del yodometano.
\end{exercise}

\begin{exercise}[$\star$]\label{exo:b2:wittig:3}
Clasifíquense como estabilizados, no estabilizados o estabilizados por un arilo: \ce{Ph3P=CHCOOEt}, \ce{Ph3P=CHCH3},
\ce{Ph3P=CHCOCH3}, \ce{Ph3P=CHC6H5}. ¿Cuál necesita butil-litio?
\end{exercise}

\begin{exercise}[$\star$]\label{exo:b2:wittig:4}
Dese el producto del fosfonoacetato de trietilo, el hidruro de sodio y el propanal, con su geometría.
\end{exercise}

\begin{exercise}[$\star\star$]\label{exo:b2:wittig:5}
Se quiere obtener metilenciclohexano. Compárese una vía de Wittig a partir de la ciclohexanona con la deshidratación
catalizada por ácido del 1-metilciclohexan-1-ol.
\end{exercise}

\begin{exercise}[$\star\star$]\label{exo:b2:wittig:6}
Desconéctese el (\textit{E})-1,2-difenileteno (estilbeno) por la \omterm{def:b2:wittig:wittig}{reacción de Wittig}. ¿Por qué la vía de Wittig
es una mala elección para el (\textit{E})-estilbeno puro, y qué se usaría en su lugar?
\end{exercise}

\begin{exercise}[$\star\star$]\label{exo:b2:wittig:7}
El propanal reacciona, por una parte, con \ce{Ph3P=CHCH3} y, por otra, con \ce{Ph3P=CHCOOEt}. Dense los
productos y su geometría principal.
\end{exercise}

\begin{exercise}[$\star\star$]\label{exo:b2:wittig:8}
Calcúlese la economía atómica de la metilenación de la ciclohexanona por \ce{Ph3P=CH2}.
\end{exercise}

\begin{exercise}[$\star\star$]\label{exo:b2:wittig:9}
Escríbase la reacción de Michaelis--Arbuzov del fosfito de trietilo con el bromoetanoato de etilo, en dos etapas:
sustitución por el fósforo y luego desalquilación por el bromuro.
\end{exercise}

\begin{exercise}[$\star\star\star$]\label{exo:b2:wittig:10}
Propóngase una síntesis del 1,6-difenilhexa-1,3,5-trieno a partir del (\textit{E})-but-2-enodial
\ce{OHC-CH=CH-CHO} mediante dos reacciones de Wittig, y dígase qué iluro y en qué cantidad.
\end{exercise}

\begin{exercise}[$\star\star\star$]\label{exo:b2:wittig:11}
La (\textit{E})-4-fenilbut-3-en-2-ona puede prepararse por una \omterm{def:b2:enolates-aldol:aldol}{condensación aldólica} o por una \omterm{def:b2:wittig:wittig}{reacción de Wittig}.
Escríbanse ambas vías y compárense.
\end{exercise}

\begin{exercise}[$\star\star\star$]\label{exo:b2:wittig:12}
Propónganse dos vías hacia el (\textit{Z})-oct-4-eno, una por la \omterm{def:b2:wittig:wittig}{reacción de Wittig} y otra a través de un alquino,
y dígase qué subproductos deja cada una.
\end{exercise}

\section{Problema: una feromona por Wittig}

\begin{problem}[{Problema de fin de semana --- la retrosíntesis de un (Z)-alqueno, el iluro no estabilizado y
su selectividad, el coste en óxido de trifenilfosfina, y la alternativa del alquino}]
\label{pb:b2:wittig:1}
Objetivo: el (\textit{Z})-tricos-9-eno, \ce{CH3(CH2)7CH=CH(CH2)12CH3}, la feromona de la mosca doméstica de la
introducción. Masas molares (\unit{g/mol}): C 12.011, H 1.008, O 15.999, P 30.974, Br 79.904.
% ledger: use:muscalure, aw:C, aw:H, aw:O, aw:P, aw:Br, ghs:nonanal, ghs:BuLi

\medskip
\noindent\textbf{Parte I --- Retrosíntesis.}
\begin{enumerate}
  \item Escríbase la fórmula molecular del objetivo y compruébese que es un alqueno.
  \item ¿Qué enlace corta una desconexión de Wittig?
  \item Dense las dos parejas posibles (compuesto carbonílico, \omterm{def:b2:wittig:ylide}{sal de fosfonio}).
  \item Para la pareja nonanal + iluro, ¿qué halogenoalcano da la \omterm{def:b2:wittig:ylide}{sal de fosfonio}?
  \item ¿Por qué un halogenoalcano primario es adecuado para esta etapa?
  \item Escríbase la formación de la \omterm{def:b2:wittig:ylide}{sal de fosfonio}.
\end{enumerate}

\medskip
\noindent\textbf{Parte II --- El iluro y la geometría.}
\begin{enumerate}[resume]
  \item ¿Qué base desprotona esta sal? ¿Por qué no serviría el hidróxido de sodio?
  \item ¿Está estabilizado el iluro? ¿Qué geometría predice eso?
  \item Escríbase el \omterm{def:b2:wittig:wittig}{oxafosfetano} formado con el nonanal.
  \item ¿Por qué debe llevarse a cabo la reacción en seco y bajo gas inerte?
  \item ¿Por qué el doble enlace queda exactamente entre los carbonos 9 y 10?
  \item ¿Qué daría en cambio la deshidratación catalizada por ácido del tricosan-9-ol?
\end{enumerate}

\medskip
\noindent\textbf{Parte III --- El óxido de trifenilfosfina.}
\begin{enumerate}[resume]
  \item Escríbase la ecuación global a partir del iluro y el nonanal.
  \item Calcúlense las masas molares del objetivo, del nonanal, de la \omterm{def:b2:wittig:ylide}{sal de fosfonio} y del
        óxido de trifenilfosfina.
  \item Calcúlese la economía atómica de la etapa de Wittig (iluro + aldehído).
  \item ¿Cómo se separa el óxido de trifenilfosfina de un hidrocarburo?
  \item ¿Por qué la formación del óxido de trifenilfosfina es la fuerza motriz?
  \item ¿Por qué la reacción HWE no conviene para este objetivo?
\end{enumerate}

\medskip
\noindent\textbf{Parte IV --- La vía del alquino.}
\begin{enumerate}[resume]
  \item ¿Qué alquino, hidrogenado, da el objetivo, y con qué catalizador?
  \item Constrúyase ese alquino a partir del dec-1-ino y un halogenoalcano: ¿qué haluro, qué base?
  \item ¿Por qué debe detenerse la hidrogenación en el alqueno, y cómo lo garantiza el catalizador?
  \item Compárense las dos vías: número de etapas, control de la geometría, subproductos.
  \item ¿Cuál es la economía atómica de la etapa de hidrogenación?
  \item Indíquese la masa de óxido de trifenilfosfina coproducida por gramo de feromona en la vía de
        Wittig.
\end{enumerate}
\end{problem}
